
Across 223 large language models evaluated on AlpacaEval 2.0, model capability predicts length-controlled preference with a partial Spearman correlation of 0.985 --- after fully accounting for response verbosity. This agreement is stronger once verbosity is removed than before: the present sample yields a bivariate Spearman $\rho$ = 0.966 (N = 223), while Dubois et al. [2024] reported a bivariate correlation of approximately 0.94; the partial correlation of 0.985 exceeds both values. Controlling for the verbosity confound reveals the signal more clearly rather than attenuating it.

This finding bears on a question at the center of LLM evaluation: does length-controlled win rate (LC\_winrate) preserve the capability ordering that human-annotated win rate (win\_rate) reflects? The LC correction in AlpacaEval 2.0 is a generalized linear model (GLM)-based debiasing procedure that regresses out response length from pairwise comparison outcomes \citep{Dubois2024LengthControlled}. Its stated goal is to remove verbosity inflation --- the tendency of AI judges to favor longer, more elaborate responses --- while leaving genuine quality differences intact. Whether it succeeds at preserving capability ordering at population scale has not been tested with confound-controlled methods.

The surface problem is well-known: verbosity biases AI-based evaluation. Zheng et al. [2023] and Shi et al. [2024] document that AI judges systematically prefer longer responses. AlpacaEval 2.0's LC correction is specifically designed to address this. A deeper problem remains unresolved, however: verbosity and capability co-vary in real model populations (VIF = 1.764 in the present data), meaning that any observed correlation between win\_rate and LC\_winrate could reflect shared verbosity rather than shared capability signal. The gap in existing work is precisely this: no study has applied partial correlation or standardized regression to disentangle the independent contributions of capability and verbosity to LC preference across a large, diverse model population.

The key observation is that the capability-LC alignment question is inherently a partial correlation problem. Once avg\_length is controlled, the capability-LC association does not weaken --- it strengthens. Verbosity acts as a mild suppressor: capable models tend to write longer responses, which partially masks the true capability-LC signal in bivariate analysis. After residualization, the underlying capability signal dominates LC prediction by a factor of approximately 4.9 (standardized $\beta$ ratio: 21.34 vs 4.37), and the negative $\beta$ for verbosity confirms the LC correction is functional.

This paper makes the following contributions:

\begin{enumerate}
\item \textbf{First population-scale partial correlation quantification}: $\rho$(win\_rate, LC\_winrate | avg\_length) = 0.985 (p = 1.69e-170, one-tailed, pre-registered directional test, N = 223), with bootstrap 95\% CI [0.976, 0.988].
\end{enumerate}

\begin{enumerate}
\item \textbf{Capability-verbosity dominance via standardized regression}: Capability contributes approximately 4.$9\times$ more to LC preference than verbosity (|$\beta _win|$ = 21.34 vs |$\beta _len|$ = 4.37); full model $R^2$ = 0.963 versus verbosity-only $R^2$ = 0.256 --- a difference of 70.7 percentage points.
\end{enumerate}

\begin{enumerate}
\item \textbf{FWL-inspired verification}: Two independent estimators converge within 1.1 percentage points ($\rho_{resid}$ = 0.9739, delta = 0.011), ruling out methodological artifact from the partial correlation procedure.
\end{enumerate}

\begin{enumerate}
\item \textbf{Quartile-level monotonicity with large effect}: KW $\epsilon ^2$ = 0.883 (H = 196.32, p = 2.63e-42), with fully monotonic LC ordering (Q1 median = 7.14 < Q2 = 14.69 < Q3 = 26.41 < Q4 = 51.62) and Dunn Q1 vs Q4 p = 1.04e-38.
\end{enumerate}

\begin{enumerate}
\item \textbf{Methodological lesson on dependent variable choice}: KW on $\Delta$ yields $\epsilon ^2$ = 0.088 versus $\epsilon ^2$ = 0.883 for LC\_winrate --- a $10\times$ attenuation from dependent variable choice alone.
\end{enumerate}

The paper is organized as follows. Section 2 reviews prior work. Section 3 describes methodology. Section 4 presents experimental design. Section 5 reports results. Section 6 discusses limitations. Section 7 concludes.

